\begin{lemma}[Kempe changes preserve properness]
Let $\alpha\ne\beta$ and let $\varphi$ be a proper partial edge-coloring.
Swapping $\alpha$ and $\beta$ on one two-color component produces another
proper partial edge-coloring.
\end{lemma}
\begin{proof}
Let $K$ be the component selected in \noderef{TwoColorComponent}, and let
$\varphi'$ be the swap from \noderef{TwoColorComponentSwap}.  Consider
distinct adjacent edge identities $e$ and $f$.  We prove that they cannot
receive the same colored value under $\varphi'$.

First suppose that either both edge-nodes belong to $K$ or neither does.
Outside $K$ the color values are unchanged.  Inside $K$ the same
transposition of the two distinct values $\alpha$ and $\beta$ is applied to
both edges, and all other values, including the uncolored value, are fixed.
In either case the operation applied to the two original values is
injective.  Equality of the two new colored values would therefore imply
$\varphi(e)=\varphi(f)$, contrary to the properness hypothesis in
\noderef{IsProperEdgeColoring}.

It remains to rule out the case in which exactly one edge-node belongs to
$K$; exchange $e$ and $f$ if necessary so that $e\in K$ and $f\notin K$.
If the two new colored values were equal, inspection of the swap shows that
the original values must be the two opposite colors $\alpha$ and $\beta$.
Indeed, a value outside this pair is fixed on both sides and would reduce to
the already-contradictory equality of original colors, while an uncolored
value cannot become colored.  Since $e$ and $f$ are adjacent, some
multigraph vertex $v$ is incident with both by \noderef{AdjacentEdges}.
The definition of \noderef{TwoColorIncidenceGraph} then gives incidence-graph
edges from the edge-node $e$ to the vertex-node $v$ and from $v$ to the
edge-node $f$.  Thus the two edge-nodes are connected and belong to the
same component.  This contradicts $e\in K$ and $f\notin K$.

All possibilities lead to distinct colors on distinct adjacent colored
edges.  Hence $\varphi'$ is proper.
\end{proof}
